\documentclass[12pt, a4paper]{article} % Dokumentenklasse

\usepackage[utf8]{inputenc} % Zeichensatz und Schrift
\usepackage[T1]{fontenc}
\usepackage{lmodern} % nice alternative is {mathpazo}
\usepackage[english]{babel}

\usepackage[left=2.5cm, right=2.5cm, top=2.5cm, bottom=2.50cm]{geometry} % Seitenformat

\usepackage{amsmath} % Mathematikzeichen
\usepackage{amsfonts}
\usepackage{amssymb}
\usepackage{mathtools}

\usepackage{booktabs} % Tabellen
\usepackage{array}
\usepackage{dcolumn}
\usepackage{tabularx}
\usepackage{threeparttable}

\usepackage{graphicx} % Grafiken
\usepackage{subfigure}

\usepackage{caption} % Paket für Tabellen- und Bildbeschriftungen
\usepackage{lipsum} % Paket für Random-Text (nur für Anschauungszwecke benötigt, kann gelöscht werden)

\usepackage{setspace} % Zeilenabstand
\onehalfspacing

\usepackage{placeins}

%% Literaturverzeichnis und Zitation % %

\usepackage[%
citestyle=authoryear-comp,%
bibstyle=JME,%
maxbibnames=99,% %maximum number of names printed in bibliography before truncation with ``et al.'' is used
maxcitenames=2,
datezeros=false,% no leading 0 if dates are printed
date=long,%
isbn=false,% show no ISBNs
natbib=true,% enable natbib-compatibility
url=false,% show no urls
backend=biber % use Biber here
]{biblatex}

% vergrößert Abstand zwischen den Quellen im Lit.verzeichnis

\setlength\bibitemsep{1.5\itemsep}

% setzt "et al." bei deutscher Sprache

\DefineBibliographyStrings{ngerman}{%
  andothers = {et\addabbrvspace al\adddot}
}


% vermeidet Witwen und Waisen

\clubpenalty10000 
\widowpenalty10000
\displaywidowpenalty=10000 



% Hier .bib-Datei einfügen

\bibliography{bspbib}

%---------------------------------------------------------------------------------------
%
%								HIER BEGINNT DAS DOKUMENT
%
%---------------------------------------------------------------------------------------


\begin{document}

% TITELSEITE

\begin{titlepage}
\begin{center}
\includegraphics[scale=0.65]{KITlogo}\\
%\Large
%Karlsruher Institut für Technologie\\
\large
Fakultät für Wirtschaftswissenschaften\\
Institut für Volkswirtschaftslehre (ECON)
\vfill{
%\\\vspace{2cm}
\Large
Bachelor Thesis in Macroeconomics\\\vspace{0.5cm}
,,US commercial banks``\\\vspace{0.5cm}
Winter semester 2019/20% \vspace{2cm}
}
\vfill{
\LARGE
% -------------------------------------------------------------------------------------
% Ihr Thema
Not yet defined
% -------------------------------------------------------------------------------------
\\ 
\vspace{0.5cm}
\normalsize
% -------------------------------------------------------------------------------------
% Nummer Ihres Themas
(Topic 1)
% -------------------------------------------------------------------------------------
}
\end{center}

\vfill{
\normalsize
% -------------------------------------------------------------------------------------
% Ihre pers\"{o}nliche Daten
\noindent Alexander Schlechter \\
Matr.-Nr. 2054108 \\
alexander.schlechter@student.kit.edu}
% -------------------------------------------------------------------------------------
\end{titlepage}
\newpage

% VERZEICHNISSE

\pagenumbering{Roman}
\tableofcontents
% \listoffigures
% \listoftables
\newpage
\pagenumbering{arabic}


% ================================================================================
% 									Hauptteil
% ================================================================================


\section{Introduction}


\section{Main part}

\subsection{General look at us banks}



\begin{figure}[hbtp]
\centering
\caption{Asset side}
\includegraphics[scale=0.3]{graphs/DescriptiveStats/OtherAnalysis_AssetDistribution_7613.png}
\end{figure}


\noindent \textit{Graph description}: The graph shows the aggregates of the main variables from the asset side of the balance sheet over time.\\

\noindent \textit{Key Observations:}
\begin{itemize}
\item loans make up the largest share of assets
\item share of trading assets have risen
\item loans and trading assets have risen more than securities in timeframe year 2000-2009
\item drop in assets in 2002 and 2008
\item share in cash has fallen
\item interest derivatives have only become existent since year 2001
\item While the share of securities fell 2005-2008, the share of cash, repos did rise
\end{itemize}

\newpage

\begin{figure}[hbtp]
\centering
\caption{Share of asset positions}
\includegraphics[scale=0.3]{graphs/DescriptiveStats/OtherAnalysis_ShareofAssets_7613.png}
\end{figure}


\begin{figure}[hbtp]
\centering
\caption{Liabilities side}
\includegraphics[scale=0.3]{graphs/DescriptiveStats/OtherAnalysis_LiabilitiesDistribution_7613}
\end{figure}

\noindent \textit{Graph description}: The graph shows the aggregates of the main variables from the liabilities side of the balance sheet over time.\\

\noindent \textit{Key Observations:}
\begin{itemize}
\item deposits as main source of funding
\end{itemize}


\begin{figure}[hbtp]
\centering
\caption{Share of liabilities positions}
\includegraphics[scale=0.3]{graphs/DescriptiveStats/OtherAnalysis_ShareofLiabilities_7613}
\end{figure}

\begin{figure}[hbtp]
\centering
\caption{Growth of assets}
\includegraphics[scale=0.3]{graphs/DescriptiveStats/OtherAnalysis_GrowthAssetsYear_7613}
\end{figure}

\noindent \textit{Graph description}: The graph shows the annual growth rate of aggregated assets of all commercial banks. Two investment banks, who did become commercial banks in 2009, are excluded. \\

\noindent \textit{Key Observations:}
\begin{itemize}
\item Three negative growth rates in year 1991, 2001, 2010
\end{itemize}

\begin{figure}[hbtp]
\centering
\caption{Growth of top 1 percent banks assets}
\includegraphics[scale=0.3]{graphs/DescriptiveStats/OtherAnalysis_GrowthAssetsYear1_7613.png}
\end{figure}

\noindent \textit{Graph description}: The graph shows the annual growth rate of aggregated assets of top top $1/%$ commercial banks. Two investment banks, who did become commercial banks in 2009, are excluded. \\

\noindent \textit{Key Observations:}
\begin{itemize}
\item More negative growth rates in 1990 and 2010
\item No negative growth in 2001

\end{itemize}

\begin{figure}[hbtp]
\centering
\caption{Growth of all banks vs top 1 percent}
\includegraphics[scale=0.3]{graphs/DescriptiveStats/OtherAnalysis_GrowthTrends_7613}
\end{figure}


\noindent \textit{Graph description}: Annual growth rate with trend for all banks and top $1/%$. \\

\noindent \textit{Key Observations:}
\begin{itemize}
\item time-frame 2000-2005: growth rate of top $1/%$ rising, growth rate of all banks falling.
\item 

\begin{figure}[hbtp]
\centering
\caption{Banks default}
\includegraphics[scale=0.3]{graphs/DescriptiveStats/NewLeverage_BanksDefaultPercentage_7613.png}
\end{figure}

\noindent \textit{Graph description}: The graph shows an estimation of how many banks have defaulted at a certain time (year,quarter). It is based on the negative equity recorded by banks. Hence, it is not exact and some banks might continue to exist in case of mergers or bailouts. Also sometimes banks are double counted, if a negative equity does not immediately result in bankruptcy. \\
\noindent \textit{Key Observations:}
\begin{itemize}
\item main defaults in years 1986-1991 and 2009-2011
\item long stable period from 1991-2008
\item In 1990 there were many more smaller banks. Smaller banks might have a higher likelihood to fail. In 1990:  $~74\%$ small banks, 2010: $~35\%$ small banks 
\end{itemize}


\begin{figure}[hbtp]
\centering
\caption{Loans}
\includegraphics[scale=0.3]{graphs/DescriptiveStats/OtherAnalysis_LoanDistribution_7613}
\end{figure}

\noindent \textit{Graph description}: It shows the share of loan types of total loans over time.\\

\noindent \textit{Key Observations:}
\begin{itemize}
\item real estate loans has largest share

\end{itemize}

\begin{figure}[hbtp]
\centering
\caption{Loans by repricing maturity}
\includegraphics[scale=0.3]{graphs/DescriptiveStats/OtherAnalysis_ShareofMaturityLoans_9713}
\end{figure}


\begin{figure}[hbtp]
\centering
\caption{Residential Loans by repricing maturity}
\includegraphics[scale=0.3]{graphs/DescriptiveStats/OtherAnalysis_ShareofMaturityResLoans_9713}
\end{figure}


\begin{figure}[hbtp]
\centering
\caption{Aggregate assets by percentiles}
\includegraphics[scale=0.3]{graphs/DescriptiveStats/OtherAnalysis_BanksPercentiles_7613}
\end{figure}

\begin{figure}[hbtp]
\centering
\caption{Top 10 banks assets vs rest}
\includegraphics[scale=0.3]{graphs/DescriptiveStats/LeverageDistribution_AssetsTop10vsRest_7813}
\end{figure}


\begin{figure}[hbtp]
\centering
\caption{Rise of top 10 banks asset share}
\includegraphics[scale=0.3]{graphs/DescriptiveStats/LeverageDistribution_shareTop10_7813}
\end{figure}

\textbf{Banks by asset size}

\begin{figure}[hbtp]
\centering
\caption{Banks count by asset size}
\includegraphics[scale=1]{graphs/MedianBank_FrequencyTable_8010}
\end{figure}

In Figure 12, we have in the left column the asset interval size and in the corresponding row the number of banks per year.\\

\textbf{Typical small/medium/large bank}

Banks are assigned three different buckets (small/medium/large) depending on asset size.\\

Small bank: $0<$ assets $<= 10^5$

Medium bank: $10^5<$ assets $<= 10^6$

Large bank: $10^6<$ assets 

\begin{figure}[hbtp]
\centering
\caption{Asset size by bank}
\includegraphics[scale=0.3]{graphs/DescriptiveStats/MedianBank_SmallvsMiddlevsBigMedianBankAssets_7613.png}
\end{figure}

\begin{figure}[hbtp]
\centering
\caption{Medium vs large bank by asset size}
\includegraphics[scale=0.3]{graphs/DescriptiveStats/MedianBank_MiddlevsBigMedianBankAssets_7613.png}
\end{figure}


\begin{figure}[hbtp]
\centering
\caption{Small bank: liability side}
\includegraphics[scale=0.3]{graphs/DescriptiveStats/MedianBank_ShareofLiabilitiesSmallBank_7613.png}
\end{figure}


\begin{figure}[hbtp]
\centering
\caption{Medium bank: liability side}
\includegraphics[scale=0.3]{graphs/DescriptiveStats/MedianBank_ShareofLiabilitiesMediumBank_7613.png}
\end{figure}


\begin{figure}[hbtp]
\centering
\caption{Large bank: liability side}
\includegraphics[scale=0.3]{graphs/DescriptiveStats/MedianBank_ShareofLiabilitiesBigBank_7613.png}
\end{figure}



\newpage

\subsection{Looking into leverage}

Throughout the analysis the definition of accounting leverage (assets/equity) is used. Equity is calculated by total assets minus total liabilities. In addition, for risk analysis banks belonging a bank holding company were aggregated. Hence, the dataset which was used contained bhcs and independent banks.

\begin{figure}[hbtp]
\centering
\caption{Average leverage over time (1976/1-2013/4)}
\includegraphics[scale=0.3]{graphs/LeverageDistribution_LeverageRatio_7613}
\end{figure}

\begin{figure}[hbtp]
\centering
\caption{Average leverage trends over time (1976/1-2013/4)}
\includegraphics[scale=0.3]{graphs/LeverageDistribution_LeverageRatioTrends_7613.png}
\end{figure}

\noindent \textit{Graph description}: The graph shows the average leverage (assets/equity) for every year over all banks. Banks with equity or assets below zero are excluded.


\noindent \textit{Key Observations:}
\begin{itemize}
\item Overall Leverage did fall over time
\item Spike in leverage in year 2008/2009
\item Leverage lowest in 2007
\item Small spike in year 1999
\item Introduction of Basel 1 in 1988 might have lead to continuously decrease in leverage
\end{itemize}

\begin{figure}[hbtp]
\centering
\caption{Boxplots (1976-2013)}
\includegraphics[scale=0.3]{graphs/Leverage/LeverageDistribution_LeverageRatioBoxplot_7613.png}
\end{figure}

\noindent \textit{Graph description}: Boxplots of all leverage ratios by banks by year. 
\\

\noindent \textit{Key Observations}:
\begin{itemize}
\item large standard deviation in year 2010 with 18.82
\item less and higher bars in 2013 indicate higher homogeneity in 2013 compared to the years before.
\end{itemize}



\iffalse
\begin{figure}[hbtp]
\centering
\caption{Leverage Top 10 vs Rest over all years}
\includegraphics[scale=0.4]{graphs/Leverage/LeverageDistribution_LeverageRatioTop10vsRest_7813}
\end{figure}


\newpage

\noindent \textit{Graph description}: Since the top 10 banks share of assets did rise up to 60\% in 2013, it is important to differentiate. The graph shows the average leverage (assets/equity) for every year quarter 4.\\ Blue Line: Top 10 Banks by assets\\
Orange Line: All banks beside the top 10\\

\noindent \textit{Key Observations}:
\begin{itemize}
\item Leverage of top 10 banks tends to be higher
\item Trend of falling leverage is similar
\end{itemize}




\begin{figure}[hbtp]
\centering
\caption{Leverage Top10 vs Rest detailed look into crisis}
\includegraphics[scale=0.4]{graphs/Leverage/LeverageDistribution1_LeverageDetailCrisisTop10vsRest_0510}
\end{figure}


\noindent \textit{Key Observations}:
\begin{itemize}
\item Top10 banks leverage peak in year 2008/3 before the rest banks leverage peak in year 2009/4 (theory of risky assets from big banks to small transfer?) 
\end{itemize}
\fi
\pagebreak

\textbf{A look into the distribution of leverage}\\



\begin{figure}[hbtp]
\centering
\caption{Distribution 1980-2013}
\includegraphics[scale=0.3]{graphs/Leverage/LeverageDistribution_LeverageRatio_8013}
\end{figure}

\noindent \textit{Graph description}: Counts are normed to 1. Only leverage ratios between 0-20 are accounted for. The small amount of banks above with leverage above 20 are seen as outliers.As you can see ~50$/%$ of banks are within a range of 8-14. Leverage ratios are always from quarter 4.
\\

\noindent \textit{Key Observations}:
\begin{itemize}
\item large standard deviation in year 2010 with 18.82
\item less and higher bars in 2013 indicate higher homogeneity in 2013 compared to the years before.
\end{itemize}

\pagebreak


\begin{figure}[hbtp]
\centering
\caption{Distribution in crisis 2003-2011}
\includegraphics[scale=0.3]{graphs/Leverage/LeverageDistribution_LeverageRatio_0311}
\end{figure}

\noindent \textit{Graph description}: Counts are normed to 1. Only leverage ratios between 0-20 are accounted for. The others are seen as outliers. As you can see ~50$/%$ of banks are within a range of 8-12. Leverage ratios are always from quarter 4.

\noindent \textit{Key Observations}:
\begin{itemize}
\item Increasing homogenity over time.
\end{itemize}


\section{Conclusion}




% ================================================================================
% 							Literaturverzeichnis
% ================================================================================

\newpage
\printbibliography[
heading=bibintoc,
title={Bibliography}
]


% ================================================================================
% 								Appendix
% ================================================================================

%\appendix
%\section{Anhang}




\end{document}